\section*{Why this book}

Computational mechanics is learned well when three things are seen together: how a mathematical formulation becomes a numerical algorithm, how that algorithm is implemented, and how its results are verified. This book is organised around that connection. It develops finite element analysis (FEA) and reduced-order modelling (ROM) through the mathematics needed to follow each step, working Python implementations, structural mechanics examples, and a systematic examination of the numerical results.

\paragraph{From mathematical formulation to working software.} Finite element analysis is usually taught through weak forms, element matrices, assembly and solution algorithms, while reduced-order modelling is introduced separately through projection, modal truncation or data-driven approximation. Implementing both shows how they connect. The book follows one path from material properties, mesh and element formulation to global equilibrium, solution and post-processing, and on to snapshot generation, basis construction and reduced-order prediction. The accompanying packages, \fea{} and \rom{}, give concrete implementations through which readers can examine these connections instead of treating the software as a black box.

\paragraph{Why write the software instead of simply using established tools?} Established finite element and model-reduction libraries serve important purposes and may be the right choice for many engineering applications. This book does not seek to replace them or claim superiority in performance or breadth. It provides a compact implementation that readers can trace, inspect, test and extend, and that needs only Python, NumPy and SciPy, with PyTorch optional. The finite element and model-reduction parts are deliberately decoupled: \rom{} works on numerical arrays and does not import \fea{}, as Chapter~\ref{ch:arch} shows and checks. This helps readers see which parts of a reduced-order workflow belong to a particular finite element solver and which can be applied to data and operators produced by other tools.

\paragraph{Structural mechanics beyond introductory examples.} The book moves from foundational formulations to the computational challenges of realistic structural analysis. Alongside element behaviour, meshes, assembly and linear analysis, it covers nonlinear equilibrium, material and geometric nonlinearity, contact, transient dynamics, computational acceleration, and reduced-order models for structural systems. Advanced topics include component-mode substructuring, nonlinear reduced models, and hyper-reduction for contact problems. Each topic is treated through its mathematical assumptions, implementation choices, computational evidence and limitations, not as a catalogue of available algorithms.

\paragraph{Verification and reproducibility.} A numerical result is useful only when its meaning and reliability can be assessed, so verification is part of the method and not a final software task. Depending on the problem, the evidence is an analytical solution, an independent computation, convergence behaviour, a consistency check, an automated test, or a reduced-model error on held-out cases. The book separates verification, whether the implementation solves the stated equations correctly, from validation, whether those equations describe physical reality, and it does not treat passing tests as proof that a program is correct for every problem. Every number in the text is produced by a script and enters the book as a macro, so the number a reader sees is the one that was measured, and readers can look into the model, solver settings, reference values, error measures and costs behind it. The book separates the cost of constructing a reduced model from its online solution, and reduced-order accuracy from speed-up. It also names the methods that are implemented but not exercised by its examples, and the cases where a result is limited by the benchmark, the implementation or the available evidence.

\paragraph{A foundation for further work.} The purpose of the book is to help readers move from using numerical tools to understanding and developing them. Graduate students can use it to connect mathematical theory with computational practice, and engineering researchers can use it to study algorithms, reproduce examples, examine implementation decisions and lay the ground for new methods. The packages are working research software and learning tools, not a claim that every included method is universally applicable or production-ready.
